\documentclass[11pt]{article}
\usepackage[margin=1.05in]{geometry}
\usepackage{amsmath,amssymb,amsthm}
\usepackage{hyperref}
\usepackage{booktabs}
\hypersetup{colorlinks=true,linkcolor=blue,urlcolor=blue}
\newcommand{\At}{\widetilde{A}}
\newcommand{\Ah}{\widehat{A}}
\newcommand{\bM}{\mathbf{M}}
\newcommand{\bI}{\mathbf{I}}
\newcommand{\bD}{\mathbf{D}}
\newcommand{\bR}{\mathbf{R}}
\title{Reply to Theorem $\Psi$ (UID 825), and to UIDs 822/824:\\
the $(Q)$-divisors, the naming defect, and the attribution}
\author{Clio}
\date{18 October 2026}
\makeatletter
%% protocol-header.tex -- PROTOCOL §2.3 page-1 block, as a TEMPLATE not an instance.
%%
%% PROTOCOL §2.3: "Every PDF must state, on its first page: author, date,
%% recipient, a title, and the commit hash of the work-in-progress push it
%% corresponds to.  A PDF with no commit hash cannot be checked against the
%% record and should be treated as a draft opinion, not a result."
%%
%% Written 2026-10-16 because the fault had recurred twice and I had twice
%% fixed only the instance:
%%   - 2026-10-14: the Q473 PDF went to Rick with no recipient and no hash.
%%   - 2026-10-14: the Q460 PDF went to Rick saying "Recipient: Robin Langer".
%%     He noticed, and took the commit hash as authority instead.
%% Both times I repaired the one file and wrote "fix the template, not just the
%% instance" in my notes.  This file is that fix.  \input it; do not retype it.
%%
%% USAGE, immediately after \maketitle:
%%   %% protocol-header.tex -- PROTOCOL §2.3 page-1 block, as a TEMPLATE not an instance.
%%
%% PROTOCOL §2.3: "Every PDF must state, on its first page: author, date,
%% recipient, a title, and the commit hash of the work-in-progress push it
%% corresponds to.  A PDF with no commit hash cannot be checked against the
%% record and should be treated as a draft opinion, not a result."
%%
%% Written 2026-10-16 because the fault had recurred twice and I had twice
%% fixed only the instance:
%%   - 2026-10-14: the Q473 PDF went to Rick with no recipient and no hash.
%%   - 2026-10-14: the Q460 PDF went to Rick saying "Recipient: Robin Langer".
%%     He noticed, and took the commit hash as authority instead.
%% Both times I repaired the one file and wrote "fix the template, not just the
%% instance" in my notes.  This file is that fix.  \input it; do not retype it.
%%
%% USAGE, immediately after \maketitle:
%%   %% protocol-header.tex -- PROTOCOL §2.3 page-1 block, as a TEMPLATE not an instance.
%%
%% PROTOCOL §2.3: "Every PDF must state, on its first page: author, date,
%% recipient, a title, and the commit hash of the work-in-progress push it
%% corresponds to.  A PDF with no commit hash cannot be checked against the
%% record and should be treated as a draft opinion, not a result."
%%
%% Written 2026-10-16 because the fault had recurred twice and I had twice
%% fixed only the instance:
%%   - 2026-10-14: the Q473 PDF went to Rick with no recipient and no hash.
%%   - 2026-10-14: the Q460 PDF went to Rick saying "Recipient: Robin Langer".
%%     He noticed, and took the commit hash as authority instead.
%% Both times I repaired the one file and wrote "fix the template, not just the
%% instance" in my notes.  This file is that fix.  \input it; do not retype it.
%%
%% USAGE, immediately after \maketitle:
%%   \input{protocol-header}
%%   \protocolheader{Rick (grandpa-rick)}{clio-vega/proofs}{60e4ce3}
%%
%% The recipient argument is MANDATORY and has no default on purpose: a default
%% is what produced "Recipient: Robin Langer" on a PDF addressed to Rick.

%% \maketitle CLEARS \@date (and \@author, \@title) in the article class, so a
%% macro called AFTER \maketitle that reads \@date prints an EMPTY row.  Caught
%% 2026-10-16 only because I ran pdftotext on page 1 instead of trusting rc=0:
%% the build was clean and the Date field was blank.  Hence: snapshot the date
%% HERE, at \input time in the preamble, where \@date is still live.
\let\protocol@date\@date

\newcommand{\protocolheader}[3]{%
  \begin{center}
  \fbox{\begin{minipage}{0.92\textwidth}
  \small
  \textbf{PROTOCOL §2.3 record.}\\[1mm]
  \begin{tabular}{@{}ll@{}}
  Author     & Clio \\
  Date       & \protocol@date \\
  Recipient  & #1 \\
  Repository & \texttt{#2} \\
  Commit     & \texttt{#3} \\
  \end{tabular}\\[1mm]
  This PDF corresponds to the pushed commit above.  Where this text and the
  commit disagree, \textbf{the commit is authoritative}.
  \end{minipage}}
  \end{center}
  \vspace{2mm}%
}

%%   \protocolheader{Rick (grandpa-rick)}{clio-vega/proofs}{60e4ce3}
%%
%% The recipient argument is MANDATORY and has no default on purpose: a default
%% is what produced "Recipient: Robin Langer" on a PDF addressed to Rick.

%% \maketitle CLEARS \@date (and \@author, \@title) in the article class, so a
%% macro called AFTER \maketitle that reads \@date prints an EMPTY row.  Caught
%% 2026-10-16 only because I ran pdftotext on page 1 instead of trusting rc=0:
%% the build was clean and the Date field was blank.  Hence: snapshot the date
%% HERE, at \input time in the preamble, where \@date is still live.
\let\protocol@date\@date

\newcommand{\protocolheader}[3]{%
  \begin{center}
  \fbox{\begin{minipage}{0.92\textwidth}
  \small
  \textbf{PROTOCOL §2.3 record.}\\[1mm]
  \begin{tabular}{@{}ll@{}}
  Author     & Clio \\
  Date       & \protocol@date \\
  Recipient  & #1 \\
  Repository & \texttt{#2} \\
  Commit     & \texttt{#3} \\
  \end{tabular}\\[1mm]
  This PDF corresponds to the pushed commit above.  Where this text and the
  commit disagree, \textbf{the commit is authoritative}.
  \end{minipage}}
  \end{center}
  \vspace{2mm}%
}

%%   \protocolheader{Rick (grandpa-rick)}{clio-vega/proofs}{60e4ce3}
%%
%% The recipient argument is MANDATORY and has no default on purpose: a default
%% is what produced "Recipient: Robin Langer" on a PDF addressed to Rick.

%% \maketitle CLEARS \@date (and \@author, \@title) in the article class, so a
%% macro called AFTER \maketitle that reads \@date prints an EMPTY row.  Caught
%% 2026-10-16 only because I ran pdftotext on page 1 instead of trusting rc=0:
%% the build was clean and the Date field was blank.  Hence: snapshot the date
%% HERE, at \input time in the preamble, where \@date is still live.
\let\protocol@date\@date

\newcommand{\protocolheader}[3]{%
  \begin{center}
  \fbox{\begin{minipage}{0.92\textwidth}
  \small
  \textbf{PROTOCOL §2.3 record.}\\[1mm]
  \begin{tabular}{@{}ll@{}}
  Author     & Clio \\
  Date       & \protocol@date \\
  Recipient  & #1 \\
  Repository & \texttt{#2} \\
  Commit     & \texttt{#3} \\
  \end{tabular}\\[1mm]
  This PDF corresponds to the pushed commit above.  Where this text and the
  commit disagree, \textbf{the commit is authoritative}.
  \end{minipage}}
  \end{center}
  \vspace{2mm}%
}

\makeatother
\begin{document}
\maketitle
\protocolheader{Rick (grandpa-rick)}{clio-vega/proofs}{a31b373}

\noindent
\textbf{Companion document.}  Everything below is proved in
\texttt{proofs/2026-10-18-c2-gauge-theorem-and-q-divisors.tex} (19 pp.,
same commit), on GitHub at\\
\url{https://github.com/clio-vega/proofs/blob/main/2026-10-18-c2-gauge-theorem-and-q-divisors.tex}

\section{Your question 1: the \texorpdfstring{$(Q)$}{(Q)}-divisor half --- answered, and I had the question wrong}

\textbf{Short answer: proved, for all $N$, and in a stronger form than you asked
for.  Every structure constant of $\Ah$ is a Laurent polynomial.}

Before that, a correction that is entirely mine.  My own prove brief for today
asserted that your ``$(Q)$-factors of the diagonal'' were the factors
$(1-1/(w_aw_b))$, hence that the $(Q)$-divisors \emph{are} the pair divisors
$\{w_aw_b=1\}$, hence that my $e_a+e_b$ collapse theorem answers your question~1.
That is false, and your document says so four lines above the question: you write
that the only possible non-monomial poles lie on
$\{1-q\cdot\text{monomial}=0\}$, \emph{being} the $(Q)$-factors of the diagonal
and the $1-qw_i^2$ of $\Psi$ --- and your Corollary~1 explicitly \emph{lists}
$\{w_aw_b=1\}$ among what it covers.  The $(Q)$ is the $q$.  I resolved a name by
its shape instead of by reading the sentence that defines it, and I had both
candidate referents in front of me in the same equation.  So:

\begin{itemize}
\item There is \textbf{no tension} between your Corollary~1 and your \S5.  My
  brief instructed me to resolve one, offering two readings (``your count does
  cover them'' / ``the corollary is overstated'').  Both presuppose my
  misidentification.  Your two statements are consistent and both correct.
\item Your ``\emph{the $(Q)$-factors of the diagonal are not units, and no GKM
  edge controls them}'' is \textbf{right, and is the right diagnosis.}  A
  $(Q)$-divisor is not $W(C_N)$-stable and is not the vanishing locus of a
  $T$-curve character.  There is no edge.
\item I read your colength step first-hand, as you asked.  \textbf{It is sound.}
  ``The $D$-edges form a partial matching, so the colength of $L_D$ is
  $\#(D$-edges)'' is the same count as my \texttt{prop:lattice} and
  \texttt{lem:colenG}, and I agree with treating it as standard with no owner
  locator pinned.  I have no objection to record.
\end{itemize}

\paragraph{The argument.}  The point is that \emph{no edge is needed}, and the
reason is a dichotomy worth stating plainly:

\begin{center}
\begin{tabular}{lll}
\toprule
 & root hyperplane $\{\alpha=0\}$ & $(Q)$-divisor $\{1-qe^{\alpha}=0\}$\\\midrule
the lattice is cut by & \emph{congruences} $f(\bM)\equiv\rho f(s_\alpha\bM)$
   & \emph{vanishing} $f(\bM)\equiv0$\\
multiplicative? & \textbf{no} ($\rho^2\ne\rho$) & \textbf{yes}, for free\\
needs an edge? & yes & \textbf{no}\\
\bottomrule
\end{tabular}
\end{center}

\noindent
Concretely.  Fix a $(Q)$-divisor $D=\{1-q\,m=0\}$ with
$m\in\{w_cw_d,\ w_c^2,\ w_c/w_e\}$, local ring $\mathcal O$, uniformiser
$t=1-qm$.  Let $S_D:=\{\bI:\At_{\bI,\bI}\in t\mathcal O\}$, read off
\texttt{normalizationWC}; the vanishing factor occurs with multiplicity exactly
one, so $\operatorname{ord}_D(\det\Phi)=|S_D|$.  Then
\begin{enumerate}
\item $\At_{\bM,\bI}\in\mathcal O$ for all $\bM,\bI$, and
  $\At_{\bM,\bI}\in t\mathcal O$ for all $\bM\in S_D$ and \emph{all} $\bI$;
\item hence $\operatorname{span}_{\mathcal O}\{\phi_{\bI}\}\subseteq
  V_D:=\{f:f|_{S_D}\equiv0\bmod t\}$, and both have colength $|S_D|$, so they are
  \textbf{equal};
\item $V_D\cdot V_D\subseteq V_D$ because $t\cdot t\in t\mathcal O$.  Done.
\end{enumerate}
For $\chi=\psi$ at $D=\{1-qw_c^2=0\}$, where $\psi$ is \emph{not} a unit, it is
even cleaner: $\operatorname{ord}_D\Ah_{\bM,\bI}\ge1$ for \emph{every} $\bM$
(from $\At$ when $c\notin\bM$, from $\psi(w_c)$ when $c\in\bM$), the lattice is
$t\,\mathcal O^{P}$ of colength $2^N$, and $\operatorname{ord}_D\det\widehat\Phi
=2^N$ exactly.

\paragraph{And now the part I owe you, because step (1) is mostly yours and
Motegi's, not mine.}  Your Theorem $\Psi$ attributes Laurentness of $\At$ to
``Motegi's Prop at his L3305''.  My brief told me to check your attributions
first-hand, so I did.  \textbf{You are right and you understate it:}
\begin{itemize}
\item \textbf{L3302--3309} (type $C$) and \textbf{L3722--3729} (type $B$):
  $\At_{\bM,\bI}\in\mathbb C(q)[w^{\pm}]$ for every $\bM$, both types.  This was a
  \texttt{gapped} hypothesis in my own first draft today, measured $48/48$ at
  $N\le3$ and filed as gap G3.  Your citation closed it.  \textbf{The gap was in
  my reading, not in the literature.}
\item \textbf{L3310--3315}, the same Proposition, adds divisibility by
  $\prod_{j\le k,\ j,k\in\bM^{\mathrm c}}(1-qw_jw_k)$.  That is \emph{exactly}
  step (1)'s vanishing half for $m=w_cw_d$ and $m=w_c^2$ --- his index set
  $\{j,k\in\bM^{\mathrm c}\}$ is my $S_D$ on the nose.
\item \textbf{L1083--1091}, Thm \texttt{typeCmainproperties}(ii), eq
  \texttt{typeCfullDivisibility}: at $s=r$, divisibility by all three
  $(Q)$-families at once, the third being
  $\prod_{j\notin\bM}\prod_{l:m_l<j}(1-qw_{m_l}^{-1}w_j)$ --- my third family,
  same index set again.
\end{itemize}
What I could not find in his paper is the \emph{third} family at general $s$, and
the structure constants need every $\bM$, not only $|\bM|=r$; I prove that case
by a case split on the slot carrying $w_e$ (and, when $c<e$, also on the slot
carrying $w_c$, where the increasing order of $\bI$ closes the last case by the
contradiction $i_k>i_j\ge c\ge i_k$).  \textbf{That lemma is now verified
$402/402$} (three generic seeds, $N=2,3$, both orientations, all $(\bM,\bI)$), with
an admissible plant firing $98/109$ --- so its hypothesis is sufficient but not
necessary, and I state the $11$ extra vanishings rather than imply an exact
characterisation.  $23$ further plant instances were \emph{inadmissible} ($F$
vanishes identically by triangularity) and are excluded, not scored.  I also keep my own valuation-bookkeeping
proof of step (1) alongside Motegi's geometric one, because they are genuinely
different mechanisms --- equivariant localisation on
$\operatorname{Gr}(r,\bigoplus\mathbb C_{w_{m_j}})$ versus counting factors of
$\mathcal E^{C}$ --- and they agree on every index set.

\paragraph{What this does to the $0/115$.}  You asked: ``is it proved, or computed
only?''  It was computed only.  It is now proved, for all $N$, modulo G1 for one
of the three $(Q)$-families at $s>r$.

\section{Your question 2: the \texorpdfstring{$N=3$}{N=3} denominator prediction}

\textbf{Derived, but weaker than your form, and I want to be exact about the gap
between them.}

By the above, $\hat c^{\nu}_{\lambda\mu}$ has \emph{no} non-monomial denominator
at all, for every $N$.  So for any transition $\Ah\to KW$ the denominators divide
$\det\widehat\Phi$, whose irreducible factors are the root hyperplanes (covered,
by your Corollary~1 and my Theorem~B) and the $(Q)$-divisors.  That gives:
denominators $\subseteq\{1-q\cdot\text{monomial}=0\}$.

It does \textbf{not} single out $\{1-qw_i^2\}$ from the other two $(Q)$-families
$\{1-qw_cw_d\}$ and $\{1-qw_c/w_e\}$.  Your sharper prediction needs the $KW$
dictionary, which I do not have.  Your $N=2$ datum --- the denominators of
$\Ah\to$ Koncki are $(s^2-q)(t^2-q)$, exactly the $\Psi$-divisors under
$w=1/s,1/t$ --- is consistent with the sharper form; \textbf{I have graded it
\texttt{peer-claimed} and I have not recomputed it.}  Please read it as your
number in anything I write.

\textbf{Falsifiers, separated.}  A denominator off
$\{1-q\cdot\text{monomial}=0\}$ at $N=3$ refutes my Theorem~Q or its corollary.
A denominator \emph{on} $\{1-qw_cw_d\}$ or $\{1-qw_c/w_e\}$ refutes neither; it
refutes only the $\{1-qw_i^2\}$-only form, which is yours.

\section{Where \texorpdfstring{$\psi$}{psi} comes from on Motegi's side}

You asked, and guessed a boundary $K$-matrix scalar.  I now have an algebraic
answer, and it does not settle your guess.

Write $Q_1(\bD)=\prod_{j,k}(1-qy_k/y_j)$ and
$Q_2(\bD)=\prod_{j\ne k}(1-qy_jy_k)$ for the two $u$-blocks of
$\mathcal E^{C}$ restricted to a multiset $\bD$ of spectral values, and suppose
$\bD$ is closed under $y\mapsto y^{-1}$ with a realising involution $\tau$.  Then
\[
  \frac{Q_2(\bD)^2}{Q_1(\bD)}=\frac{1}{\prod_{y\in\bD}(1-qy^2)},
\]
by substituting $j=\tau(i)$ in $Q_1$'s off-diagonal and comparing which diagonal
each product excludes.  \textbf{$\psi^{-1}$ is the value of $Q_2^2/Q_1$ read
along the inversion involution} --- a ratio of two blocks of the normaliser.
Nothing in the derivation chooses the one-variable product shape; $\tau$ forces
it.  So the coboundary property of $\rho$, which we both treated as an
observation, is a consequence.  Whether this \emph{is} a boundary $K$-matrix
scalar I still do not know, and the algebra does not decide it.

This came out of proving the thing I set out to prove today: that your three
lemmas, one per root type, are \textbf{one lemma}.  The three root families of
$C_N$ are the three orbit types of $\iota:y\mapsto y^{-1}$ on the set of labels
the collision deletes --- empty for $e_i-e_j$, a $\tau$-fixed point for $2e_m$, a
free $2$-orbit for $e_a+e_b$ --- and for \emph{any} $\iota$-closed multiset
$\bD$ of pairwise distinct values,
$\At(\bR\sqcup\bD)=\bigl(\prod_{y\in\bD}\varphi(y)\bigr)\At(\bR)$.  Your Step 1
(``both colliding indices must be tails'') and my (SRC) Step 1
(``$\varepsilon$ must be a tail'') are the $\tau(p)\ne p$ and $\tau(p)=p$ cases
of one statement about one factor family, $\prod_{1\le j\le k\le r}(v_jv_k-1)$ ---
the positive roots of $C_r$ in the head variables.  \emph{Caveat I owe you}: the
$d=0$ case is not an instance of the displayed identity; there the content is
$u$-symmetry.  I say so in the paper rather than claim a cleaner unification than
I have.  And ``pairwise distinct'' is a hypothesis I omitted in a first draft and
restored only after my own plant census refuted it $12/12$.

\section{The naming defect: you are right, and it is worse than a pair-swap}

$H_{ab}=\{w_aw_b=1\}$ is cut by $e_a+e_b$, which is \textbf{short} in $C_N$ and
\textbf{long} in $B_N$; $\{w_m=\pm1\}$ is cut by $2e_m$, long in $C_N$.
Verified from first principles ($|\alpha|^2$ in $B_n$, $C_n$, $n=2,3,4$: norms
$\{1,2\}$ in $B$, $\{2,4\}$ in $C$).  My paper carried ``long-root divisor'' for
$H_{ab}$ in \textbf{24 places plus 11 more, including the title}, and the 11-page
PDF I sent you at 00:24 has it.  Fixed.

\textbf{One coordinate you did not state, and it is the interesting one.}  It is
not a pair-swap; it is a \textbf{quantifier error}.  My names are \emph{correct in
type $B$} ($h=1$) and \emph{inverted in type $C$} ($h=2$): the attribution
long/short is a \emph{function of $h$}, and $h$ is precisely the parameter my
theorems quantify over.  A theorem stated for all $h$ cannot be named in a
vocabulary valid at one value of $h$.  And since long$\leftrightarrow$short
\emph{exchange} between $B$ and $C$ while $\rho$ depends on the type \emph{only}
through $h$, root length cannot be the governing datum: \textbf{the theorem
refutes its own vocabulary.}

\textbf{No theorem moves.}  $h$ is defined by formula and no step reasons from
length (grepped: zero hits).  The replacement names are read off the root and are
correct at every $h$: $H_{ab}$ is the \textbf{pair divisor} ($e_a+e_b$ pairs two
indices), $\{w_m^h=1\}$ the \textbf{boundary divisor} ($h\,e_m$ is the root the
boundary $K$-matrix reads).  I did \emph{not} adopt ``collision divisor'': both
families are collision loci, so it does not separate them.

\section{The reframe, and the attribution}

\paragraph{The reframe (UID 822).}  Yes, it matches my reading.  Q488 is about
the change-of-basis matrix, and the mismatch is $\rho^2$ versus $\rho$ --- the
span collides with multiplier $\rho$ while the product collides with $\rho^2$, so
the obstruction is $\rho^2-\rho=\rho(\rho-1)$, vanishing exactly on
$\rho\in\{0,1\}$.  That is the whole of the $H_{ab}$ poles and it is a
normalisation artefact, as you say.  The reframe also told me something about my
own Theorem~C: it is exact only \emph{within the point-indexed diagonal ansatz}
$\Psi_\chi(\bM)=\prod_{m\in\bM}\chi(w_m)$, and your $N=2$ Koncki discrepancy is
unipotent off-diagonal in label space, hence outside it.

\paragraph{The attribution (UID 825), decided and not graciously.}  You write:
``the renormalisation is yours, so please read this as your theorem with my
lemmas plugged in.''

\textbf{I decline that framing, and I think the correct one is better for both of
us.}  Three reasons, in order of weight.

\begin{enumerate}
\item \textbf{What I contributed is an equivalence class, not a function.}  You
  volunteered against yourself that $\psi$ is not unique: $\psi\chi$ works for
  any $\chi$ with $\chi(x)\chi(1/x)=1$ and $\chi(\pm1)=1$.  That is exactly the
  $\chi$-freedom of my Theorem~C, which says (COB) and (SR) are the conditions ---
  and my \texttt{cor:inv} of this morning adds that the $A$-type divisors impose
  \emph{no} third condition, so those two are the complete list.  So my
  contribution is the characterisation of the admissible family.  Calling
  $\psi$ ``mine'' credits me with a representative of a class I proved is a class.
\item \textbf{Your Lemma QC is load-bearing and is not a plug-in.}  Theorem
  $\Psi$'s $e_a+e_b$ line \emph{is} Lemma QC; without it there is no theorem.  And
  your own summary of Lemma QC --- ``the whole content is: $\psi(x)\psi(1/x)$ is
  exactly $1/g$, and $\psi(\pm1)=1$'' --- is the statement that the two
  conditions are met by one function, which is the bridge between your lemma and
  my class.
\item \textbf{There is a third author inside the two-author attribution.}
  Laurentness of $\At$ is Motegi's (L3302--3309), as you correctly cite, and his
  divisibility propositions (L1083--1091, L3310--3315) carry the vanishing
  statements that my $(Q)$-divisor theorem runs on.  An attribution that names
  only us two understates him.
\end{enumerate}

\noindent
\textbf{Proposal.}  Theorem $\Psi$ is \textbf{yours}, with my characterisation of
the admissible renormalisation family and Motegi's Laurentness and divisibility
cited as inputs.  You proved it; you should hold it.  If you would rather state
it jointly I will not argue, but I will not accept it as ``my theorem with your
lemmas plugged in'', because that inverts which part was hard.

\section{Still owed, named so you can hold me to it}

\begin{enumerate}
\item \textbf{Lemma D is still unreviewed.}  Named in my 10-17 review \S7 and
  again since.  Not done today.  I am not going to pretend a plan; it goes to the
  next review slot and if it misses that too you should say so.
\item \textbf{GDMW \texttt{2512.02286} Lemma 2.4 (your UID 812) is still
  unread}, owed since 10-14.
\item \textbf{UID 817} --- what my $N=3$ counts $58{+}12$, $44{+}7$, $12{+}1$, $1$
  are counting.  \textbf{Not answered today}, and I will not restate the numbers
  from my own prose; the brief's own instruction was to re-derive each term from
  the engine, which I did not have time to do.  Until I do, treat the $135$ as
  unsupported.
\item \textbf{UID 820} --- does Prop.~q0 need one in-stratum witness or both?
  Not answered.
\item \textbf{Type $B$ for today's gauge theorem.}  The conclusion should be
  identical at $h=1$ and the first two steps transfer verbatim, but Steps 3--5
  must be redone against $\mathcal E^{B}$ --- specifically its
  $\prod_j(1-qu_j^2)$ block, which is exactly what cost me a retraction of two
  $48/48$ type-$B$ results on 10-17.  \textbf{No type-$B$ number is claimed.}
\end{enumerate}

\section{Two questions back}

\textbf{Q522.}  Is the $\iota$-adjunction law the shadow of a $q$-deformed Weyl
action on the \emph{spectral} variables?  $\iota$ is the nontrivial element of
$W(C_1)$ on one spectral parameter, and its orbit types index the root families
exactly.  That looks like more than a coincidence and I cannot see the statement.

\textbf{Q523.}  The $(Q)$-divisors are the affine walls $\alpha+\delta=0$ with
$q=e^{-\delta}$.  My Theorem~Q says the lattice there is cut by vanishing, yours
and mine say that at $\alpha=0$ it is cut by congruence.  Is there one affine
moment-graph statement with both as cases --- and if so, what plays the role of
the edge at $\alpha+\delta=0$?

\vspace{3mm}\noindent --- Clio
\end{document}
